
\documentclass{article}
\usepackage{colortbl}
\usepackage{makecell}
\usepackage{multirow}
\usepackage{supertabular}

\begin{document}

\newcounter{utterance}

\centering \large Interaction Transcript for game `cladder', experiment `full\_v1.5\_default', episode 9350 with score{-}Qwen3.5{-}9B{-}sp.
\vspace{24pt}

{ \footnotesize  \setcounter{utterance}{1}
\setlength{\tabcolsep}{0pt}
\begin{supertabular}{c@{$\;$}|p{.15\linewidth}@{}p{.15\linewidth}p{.15\linewidth}p{.15\linewidth}p{.15\linewidth}p{.15\linewidth}}
    \# & \multicolumn{2}{c}{Player} && \multicolumn{2}{c}{Game Master} \\
    \hline

    \theutterance \stepcounter{utterance}  
    & & & \multicolumn{4}{p{0.6\linewidth}}{
        \cellcolor[rgb]{0.9,0.9,0.9}{
            \makecell[{{p{\linewidth}}}]{
                \texttt{\tiny{[P1$\langle$GM]}}
                \texttt{You are an expert in causal inference. The following question is not a typical commonsense query, but rather a meticulously designed question created by a professor specializing in causal inference, intended to assess the students' mastery of the course content.} \\
\texttt{Imagine a self{-}contained, hypothetical world with only the following conditions, and without any unmentioned factors or causal relationships: Jyka has a direct effect on zory and zuph. Zuph has a direct effect on glimx. Zory has a direct effect on glimx. The overall probability of zuph is 34\%. The probability of not zuph and glimx is 64\%. The probability of zuph and glimx is 28\%. Is the chance of glimx smaller when observing zuph?} \\
\texttt{Start your answer with 'Yes' or 'No', followed by additional reasoning or evidence to support your explanation.} \\
            }
        }
    }
     \\ \\

    \theutterance \stepcounter{utterance}  
    & \multicolumn{4}{p{0.6\linewidth}}{
        \cellcolor[rgb]{0.561,0.737,0.561}{
            \makecell[{{p{\linewidth}}}]{
                \texttt{\tiny{[P1$\rangle$GM]}}
                \texttt{No} \\
\\ 
\texttt{**Reasoning:**} \\
\\ 
\texttt{To determine if the chance of `glimx` is smaller when observing `zuph`, we must compare the conditional probability \$P(\textbackslash{}text\{glimx\} \textbackslash{}mid \textbackslash{}text\{zuph\})\$ with the marginal probability \$P(\textbackslash{}text\{glimx\})\$.} \\
\\ 
\texttt{1.  **Calculate the marginal probability of `glimx` (\$P(\textbackslash{}text\{glimx\})\$):**} \\
\texttt{    We are given the probability of the joint event "not zuph and glimx" and "zuph and glimx".} \\
\texttt{    *   \$P(\textbackslash{}neg \textbackslash{}text\{zuph\} \textbackslash{}cap \textbackslash{}text\{glimx\}) = 64\textbackslash{}\%\$} \\
\texttt{    *   \$P(\textbackslash{}text\{zuph\} \textbackslash{}cap \textbackslash{}text\{glimx\}) = 28\textbackslash{}\%\$} \\
\\ 
\texttt{    Since these two events cover all possibilities for `glimx` (either `zuph` is true or false), we sum them to find the total probability of `glimx`:} \\
\texttt{    \$\$P(\textbackslash{}text\{glimx\}) = P(\textbackslash{}neg \textbackslash{}text\{zuph\} \textbackslash{}cap \textbackslash{}text\{glimx\}) + P(\textbackslash{}text\{zuph\} \textbackslash{}cap \textbackslash{}text\{glimx\})\$\$} \\
\texttt{    \$\$P(\textbackslash{}text\{glimx\}) = 64\textbackslash{}\% + 28\textbackslash{}\% =} \\
            }
        }
    }
    & & \\ \\

    \theutterance \stepcounter{utterance}  
    & & & \multicolumn{2}{p{0.3\linewidth}}{
        \cellcolor[rgb]{0.9,0.9,0.9}{
            \makecell[{{p{\linewidth}}}]{
                \texttt{\tiny{[GM$|$GM]}}
                \texttt{No} \\
            }
        }
    }
    & & \\ \\

    \theutterance \stepcounter{utterance}  
    & & & \multicolumn{2}{p{0.3\linewidth}}{
        \cellcolor[rgb]{0.9,0.9,0.9}{
            \makecell[{{p{\linewidth}}}]{
                \texttt{\tiny{[GM$|$GM]}}
                \texttt{yes} \\
            }
        }
    }
    & & \\ \\

    \theutterance \stepcounter{utterance}  
    & & & \multicolumn{2}{p{0.3\linewidth}}{
        \cellcolor[rgb]{0.9,0.9,0.9}{
            \makecell[{{p{\linewidth}}}]{
                \texttt{\tiny{[GM$|$GM]}}
                \texttt{game\_result = LOSE} \\
            }
        }
    }
    & & \\ \\

\end{supertabular}
}

\end{document}
